% Paper 28: QED on S^3 — Spectral Geometry of Perturbative Transcendentals
% GeoVac Project, April 2026
% v1.0

\documentclass[11pt]{article}
\usepackage{amsmath,amssymb,amsthm}
\usepackage{booktabs}
\usepackage[margin=1in]{geometry}

\newtheorem{theorem}{Theorem}
\newtheorem{proposition}{Proposition}
\newtheorem{corollary}{Corollary}

\title{QED on $S^3$: Spectral Geometry of Perturbative Transcendentals}
\author{GeoVac Project}
\date{April 2026}

\begin{document}
\maketitle

\begin{abstract}
We compute perturbative QED spectral sums on the three-sphere $S^3$
using the Dirac and Hodge-1 (photon) spectra with SO(4) vertex selection
rules.  Three structural results emerge.
(1)~A \emph{three-axis transcendental taxonomy}: operator order
(1st vs 2nd) determines whether Riemann even-zeta or odd-zeta appears;
bundle type (scalar vs spinor) modulates coefficients; and diagram
topology (vertex parity) introduces Dirichlet $L$-function characters,
exposing Catalan's constant $G = \beta(2)$ and Dirichlet $\beta(4)$.
(2)~An \emph{irreducible constant} $S_{\min}$ at two loops: a depth-2
multiple Hurwitz zeta value at half-integer shift that lies outside all
standard transcendental bases (15 PSLQ attempts across 47 basis elements
fail at 150-digit precision).
(3)~A \emph{depth-$k$ tower}: at three loops, a new irreducible constant
appears that is independent of $S_{\min}$ and all lower-loop constants,
confirming that each loop order produces genuinely new transcendental
content through nested vertex-parity filtration.
The transcendental class of each spectral sum is $R$-independent
(verified via Weyl's law), establishing these as structural properties
of the $S^3$ geometry rather than finite-radius artifacts.
\end{abstract}

% ===================================================================
\section{Introduction}
% ===================================================================

The GeoVac framework~\cite{paper7,paper14} computes quantum-mechanical
observables from the spectral data of graph Laplacians on compact
manifolds.  On the three-sphere $S^3$, the Dirac operator has a discrete
spectrum with known eigenvalues and degeneracies~\cite{camporesi1996},
and the framework produces structurally sparse qubit
Hamiltonians~\cite{paper14} with $O(Q^{2.5})$ Pauli scaling.

Paper~18~\cite{paper18} introduced an exchange-constant taxonomy that
classifies where transcendental numbers enter quantum-mechanical
computations: intrinsic (from angular structure), calibration (from the
$S^3$ projection), and embedding (from the electron--electron cusp).
The present paper extends this taxonomy to perturbative QED by computing
spectral sums that probe the interaction structure of quantum
electrodynamics on~$S^3$.

The central finding is that the $S^3$ spectral geometry produces the
\emph{same} transcendental constants that appear in flat-space
perturbative QED---$\zeta(3)$, $\zeta(5)$, Catalan's $G$, Dirichlet
$\beta(4)$---through a purely geometric mechanism (operator order plus
vertex topology), together with \emph{new} irreducible constants at each
loop order that have no known closed form.

% ===================================================================
\section{Spectral Data on $S^3$}
\label{sec:spectra}
% ===================================================================

\subsection{Dirac spectrum}

On the unit three-sphere, the Dirac operator has eigenvalues and
degeneracies~\cite{camporesi1996}
\begin{equation}
  |\lambda_n| = n + \tfrac{3}{2}, \qquad
  g_n = 2(n+1)(n+2), \qquad n = 0, 1, 2, \ldots
  \label{eq:dirac_spectrum}
\end{equation}
The Dirac Dirichlet series is
$D(s) = \sum_{n=0}^{\infty} g_n / |\lambda_n|^s$.
Via Hurwitz zeta,
\begin{equation}
  D(s) = 2\,\zeta(s{-}2, \tfrac{3}{2})
       - \tfrac{1}{2}\,\zeta(s, \tfrac{3}{2}).
  \label{eq:dirac_hurwitz}
\end{equation}

\subsection{Hodge-1 (photon) spectrum}

The transverse vector Laplacian (Hodge-1 operator) on $S^3$ has
eigenvalues and degeneracies
\begin{equation}
  \mu_q = q(q+2), \qquad d_q^T = q(q+2), \qquad q = 1, 2, 3, \ldots
  \label{eq:hodge1_spectrum}
\end{equation}
The Ricci shift $+2$ relative to the scalar Laplacian eigenvalues
$\ell(\ell+2)$ follows from the Bochner--Weitzenb\"ock identity
$\Delta_1 = \nabla^*\nabla + \mathrm{Ric}$ on $S^3$
(Ricci scalar $R = 6$ on unit $S^3$, giving $\mathrm{Ric} = 2g$).
There are no zero modes ($\beta_1 = 0$ for $S^3$).

\subsection{Vertex selection rule}

The QED vertex on $S^3$ couples spinor harmonics at levels $n_1, n_2$
to a vector harmonic at level $q$ with the SO(4) selection rule:
\begin{equation}
  |n_1 - n_2| \le q \le n_1 + n_2, \qquad
  n_1 + n_2 + q \;\text{is odd}.
  \label{eq:vertex_rule}
\end{equation}
The parity constraint arises from the $\gamma^\mu$ coupling, which
flips the chirality representation of $\mathrm{SO}(4)
= \mathrm{SU}(2)_L \times \mathrm{SU}(2)_R$.

% ===================================================================
\section{One-Loop Structure}
\label{sec:one_loop}
% ===================================================================

The one-loop QED effective action involves traces of $D^2$ (the squared
Dirac operator).  The Seeley--DeWitt heat kernel expansion on unit $S^3$
gives coefficients $a_0 = a_1 = \sqrt{\pi}$, $a_2 = \sqrt{\pi}/8$
(computed from the Dirac spectral data, exact in \texttt{sympy}).

The vacuum polarization coefficient and QED beta function follow:
\begin{equation}
  \Pi = \frac{1}{48\pi^2}, \qquad
  \beta(\alpha) = \frac{2\alpha^2}{3\pi}.
  \label{eq:one_loop}
\end{equation}
Both reproduce the standard flat-space one-loop QED results from $S^3$
spectral data alone.

\begin{theorem}[T9, Tier~3]
The spectral zeta function of $D^2$ satisfies
$\zeta_{D^2}(s) = 2^{2s-1}[\lambda(2s{-}2) - \lambda(2s)]$
where $\lambda(2k) = (1 - 2^{-2k})\zeta_R(2k)$.
At every integer~$s$, $\zeta_{D^2}(s)$ is a two-term polynomial in
$\pi^2$ with rational coefficients.  No odd-zeta content appears at
any~$s$.
\end{theorem}

\noindent
The mechanism: squaring maps $\mathrm{odd}^{-s}$ to
$\mathrm{odd}^{-2s}$, always producing $\pi^{\mathrm{even}}$ via
Bernoulli numbers.  Operator order (2nd vs 1st) is the primary
transcendental discriminant.

% ===================================================================
\section{Two-Loop Structure: The Hurwitz Discriminant}
\label{sec:two_loop}
% ===================================================================

At two loops, the effective action probes the \emph{first-order} Dirac
eigenvalue structure through products of propagators $1/|\lambda_n|^s$.
The key result is the even/odd $s$-parity discriminant: using the
identity $\zeta(s, 3/2) = (2^s - 1)\zeta_R(s) - 2^s$ and
Eq.~\eqref{eq:dirac_hurwitz},

\begin{equation}
  D(s) = 2(2^{s-2} - 1)\,\zeta_R(s{-}2)
       - \tfrac{1}{2}(2^s - 1)\,\zeta_R(s).
  \label{eq:hurwitz_discriminant}
\end{equation}

At even $s$, both $\zeta_R(s{-}2)$ and $\zeta_R(s)$ are
$\pi^{\mathrm{even}}$, so $D(s_{\mathrm{even}})$ is purely
$\pi^{\mathrm{even}}$.  At odd~$s$, $\zeta_R(s{-}2)$ and $\zeta_R(s)$
are odd-zeta values, producing the first appearance of $\zeta(3)$,
$\zeta(5)$, etc.

\begin{table}[h]
\centering
\caption{Dirac Dirichlet series $D(s)$: PSLQ identifications at 80-digit precision.}
\label{tab:D_values}
\begin{tabular}{@{}cll@{}}
\toprule
$s$ & $D(s)$ & Class \\
\midrule
4 & $\pi^2 - \pi^4/12$ & $\pi^{\mathrm{even}}$ \\
5 & $14\,\zeta(3) - \tfrac{31}{2}\,\zeta(5)$ & odd-zeta \\
6 & $\pi^4/3 - \pi^6/30$ & $\pi^{\mathrm{even}}$ \\
7 & $62\,\zeta(5) - \tfrac{127}{2}\,\zeta(7)$ & odd-zeta \\
8 & $2\pi^6/15 - 17\pi^8/1260$ & $\pi^{\mathrm{even}}$ \\
\bottomrule
\end{tabular}
\end{table}

\noindent
The appearance of $\zeta(3)$ at $s = 5$ is the $S^3$ analog of the
$\zeta(3)$ that enters flat-space two-loop QED through the
Schwinger--Rosner vertex correction~\cite{rosner1967}.

% ===================================================================
\section{Vertex Topology and Dirichlet $L$-Values}
\label{sec:vertex}
% ===================================================================

The two-loop sunset diagram on $S^3$ involves summing over electron
pairs $(n_1, n_2)$ with the vertex parity constraint
$n_1 + n_2 + q$ odd from Eq.~\eqref{eq:vertex_rule}.  Splitting the
Dirac Dirichlet series into even-$n$ and odd-$n$ sub-sums via
Hurwitz zeta at quarter-integer shifts $\zeta(s, 3/4)$ and
$\zeta(s, 5/4)$:

\begin{align}
  D_{\mathrm{even}}(4)
    &= \tfrac{\pi^2}{2} - \tfrac{\pi^4}{24}
       - 4G + 4\beta(4),
  \label{eq:D_even} \\
  D_{\mathrm{odd}}(4)
    &= \tfrac{\pi^2}{2} - \tfrac{\pi^4}{24}
       + 4G - 4\beta(4),
  \label{eq:D_odd} \\
  D(4) &= \pi^2 - \tfrac{\pi^4}{12}
  \qquad\text{(Catalan and $\beta(4)$ cancel)}.
  \label{eq:D_full}
\end{align}

\noindent
Here $G = \beta(2) = \sum_{n=0}^{\infty} (-1)^n/(2n+1)^2$ is
Catalan's constant and $\beta(s) = L(s, \chi_{-4})$ is the Dirichlet
beta function.  The mechanism: vertex parity acts as the Dirichlet
character $\chi_{-4}$ on the mode sum, and the quarter-integer Hurwitz
shifts expose $L(s, \chi_{-4})$ in the even/odd difference.

Equations~\eqref{eq:D_even}--\eqref{eq:D_odd} are verified by PSLQ at
80-digit precision with basis $\{1, \pi^2, \pi^4, G, \beta(4)\}$.

\subsection{Three-axis taxonomy}

The transcendental content of any QED spectral sum on $S^3$ is
classified by three independent axes:
\begin{enumerate}
\item \textbf{Operator order.}
      Second-order ($D^2$) $\to$ $\pi^{\mathrm{even}}$ (T9 theorem);
      first-order ($|D|$) $\to$ odd-zeta.
\item \textbf{Bundle type.}
      Scalar vs spinor modulates coefficients but not the
      transcendental class.
\item \textbf{Diagram topology.}
      Vertex parity introduces Dirichlet $L$-function characters
      ($G$, $\beta(4)$) that are absent from unrestricted sums.
\end{enumerate}

% ===================================================================
\section{Irreducible Constants and the Depth-$k$ Tower}
\label{sec:irreducible}
% ===================================================================

\subsection{Two-loop irreducible: $S_{\min}$}

The SO(4) Clebsch--Gordan channel count $W(n_1, n_2, q) \in \{0,1,2\}$
weights the sunset sum non-uniformly, producing the CG-weighted sum
\begin{equation}
  S_{\min} = \sum_{k=1}^{\infty} T(k)^2, \qquad
  T(k) = 2\,\zeta(2, k + \tfrac{3}{2})
       - \tfrac{1}{2}\,\zeta(4, k + \tfrac{3}{2}).
  \label{eq:S_min}
\end{equation}
This is a depth-2 multiple Hurwitz zeta value at half-integer shift
$a = 3/2$, intrinsic to the Dirac spectrum on $S^3$.

Computed at 150-digit precision with 10,000 terms and three-term tail
correction:
\begin{equation}
  S_{\min} = 2.47953\,69980\,27334\ldots
  \label{eq:S_min_value}
\end{equation}
Fifteen PSLQ attempts across 47 basis elements---including
polylogarithms $\mathrm{Li}_k(1/2)$, Euler sums, double Hurwitz
$\zeta_2(s_1,s_2; 3/2)$, polygamma $\psi^{(n)}(3/2)$, Dirichlet
$\beta(3) = \pi^3/32$, and all cross-products---return no integer
relation.  $S_{\min}$ is irreducible in every tested basis.

\subsection{Three-loop: a new irreducible constant}

The three-loop iterated sunset has chain topology:
three electron lines $(n_1, n_2, n_3)$ connected by two photon lines
$(q_1, q_2)$, with vertex selection at each junction.  The sum
factorizes as
\begin{equation}
  S^{(3)} = \sum_{n_2=0}^{\infty}
    \frac{g_{n_2}}{|\lambda_{n_2}|^a}\,V(n_2)^2, \qquad
  V(n_2) = \sum_{\substack{n, q \\ \text{allowed}}}
    \frac{g_n\,d_q^T}{|\lambda_n|^a\,\mu_q^p},
  \label{eq:three_loop_factorized}
\end{equation}
where $V(n_2) \sim 1.66 \log n_2$ (logarithmic growth, verified
numerically).  This factorization reduces the computational complexity
from $O(N^5)$ to~$O(N^3)$.

At electron exponent $a = 8$ and photon exponent $p = 2$ (chosen for
fast convergence, $\sim$20 digits at $N = 50$), PSLQ with an
18-element basis \emph{including $S_{\min}$ and all its cross-products}
with $\pi^2$, $\zeta(3)$, $G$, $D(8)$, etc.\ returns no relation.

\begin{proposition}[Depth-$k$ tower]
At $k$ loops on $S^3$ with vertex restrictions and CG weights:
\begin{itemize}
\item $k = 1$: $\pi^{\mathrm{even}}$ (T9 theorem, proven).
\item $k = 2$: $S_{\min}$ (irreducible, 15 PSLQ failures, 47 elements).
\item $k = 3$: new constant $S^{(3)}$ (irreducible, independent of
      $S_{\min}$, PSLQ with $S_{\min}$ in basis fails).
\end{itemize}
Each loop order produces a genuinely new transcendental constant through
nested vertex-parity filtration.
\end{proposition}

The mechanism: each vertex coupling creates one level of Hurwitz nesting.
One vertex produces depth-2 (the min-weight sum in $S_{\min}$).
Two vertices produce depth-3 (nested min-weights in $S^{(3)}$).
The prediction extends: $k$ vertices should produce depth-$(k+1)$
multiple Hurwitz zeta values.

% ===================================================================
\section{Flat-Space Consistency}
\label{sec:flat_space}
% ===================================================================

On $S^3$ of radius $R$, the Dirac eigenvalues scale as
$|\lambda_n(R)| = (n + 3/2)/R$ and the Dirichlet series becomes
$D(s, R) = R^s \cdot D(s, 1)$.  The ratio $D(s,R)/R^s$ is
$R$-independent \emph{exactly}, confirming that the even/odd
$s$-parity discriminant is structural, not a finite-$R$ artifact.

Weyl's law is verified: the cumulative Dirac mode count
$N(n_{\max}) = 2(n_{\max}+1)(n_{\max}+2)(n_{\max}+3)/3$ matches the
Weyl prediction $(2/3)R^3 p_{\max}^3$ with $O(1/n_{\max})$
convergence.

The \emph{same} transcendental constants---$\zeta(3)$, Catalan $G$,
Dirichlet $\beta(4)$---appear in both the $S^3$ spectral sums and
flat-space two-loop QED~\cite{rosner1967,laporta1996}.
Coefficient matching for the full Petermann--Sommerfield anomalous
magnetic moment $a_2 = -197/144 + \pi^2/12 + 3\zeta(3)/4 -
\pi^2\ln 2/2$~\cite{petermann1957,sommerfield1957} requires matching
specific diagram topologies and renormalization subtractions, which is
beyond the scope of spectral sums.

% ===================================================================
\section{Discussion}
\label{sec:discussion}
% ===================================================================

The three-axis taxonomy extends Paper~18's exchange-constant
classification~\cite{paper18} from free propagators to interaction
vertices.  The key structural insight is that \emph{diagram topology}
is an independent source of transcendental content: unrestricted sums
produce only the constants dictated by operator order, while vertex
restrictions expose Dirichlet $L$-function characters and, at higher
depth, irreducible multiple Hurwitz zeta values.

The depth-$k$ tower has a natural interpretation in the framework of
Zagier's motivic weight conjecture~\cite{zagier1994}: each loop order
increases the motivic weight by the vertex coupling, and the algebra
of multiple zeta values provides the natural receptacle for the
resulting transcendentals.  The half-integer shift $a = 3/2$ from the
Dirac spectrum on $S^3$ produces Hurwitz (rather than Riemann)
multiple zeta values, and the failure of standard bases suggests that
the $S^3$ spectral geometry accesses a sector of the MZV algebra that
is not reached by flat-space Feynman integrals at the same loop order.

\subsection{Open questions}

\begin{enumerate}
\item Can $S_{\min}$ be identified with a known multiple Hurwitz zeta
      value in the Borwein--Bailey--Broadhurst
      compendium~\cite{bailey2001}?
\item Does the depth-$k$ tower stabilize (i.e., do all constants
      beyond some $k_0$ reduce to products of lower-depth constants)?
\item Is there a Picard--Fuchs interpretation of $S_{\min}$, analogous
      to Ap\'ery's recurrence for $\zeta(3)$~\cite{apery1979}?
\end{enumerate}

% ===================================================================
% References
% ===================================================================

\begin{thebibliography}{99}

\bibitem{paper7}
GeoVac Project, ``The Dimensionless Vacuum: $S^3$ Conformal Equivalence
and 18 Symbolic Proofs,'' Paper~7 (2025).

\bibitem{paper14}
GeoVac Project, ``Structurally Sparse Qubit Hamiltonians from Natural
Geometry,'' Paper~14 (2025).

\bibitem{paper18}
GeoVac Project, ``Spectral-Geometric Exchange Constants,'' Paper~18
(2026).

\bibitem{camporesi1996}
R.~Camporesi and A.~Higuchi, ``On the eigenfunctions of the Dirac
operator on spheres and real hyperbolic spaces,''
\textit{J.~Geom.~Phys.}\ \textbf{20}, 1--18 (1996).

\bibitem{rosner1967}
J.~L.~Rosner, ``Higher-order contributions to the divergent part of
$Z_3$ in a model quantum electrodynamics,''
\textit{Ann.~Phys.}\ \textbf{44}, 11--34 (1967).

\bibitem{laporta1996}
S.~Laporta and E.~Remiddi, ``The analytical value of the electron
$(g-2)$ at order $\alpha^3$ in QED,''
\textit{Phys.~Lett.~B}\ \textbf{379}, 283--291 (1996).

\bibitem{petermann1957}
A.~Petermann, ``Fourth-order magnetic moment of the electron,''
\textit{Helv.~Phys.~Acta}\ \textbf{30}, 407--408 (1957).

\bibitem{sommerfield1957}
C.~M.~Sommerfield, ``Magnetic dipole moment of the electron,''
\textit{Phys.~Rev.}\ \textbf{107}, 328--329 (1957).

\bibitem{zagier1994}
D.~Zagier, ``Values of zeta functions and their applications,''
in \textit{First European Congress of Mathematics}, Vol.~II,
Birkh\"auser, 497--512 (1994).

\bibitem{bailey2001}
D.~H.~Bailey and D.~J.~Broadhurst, ``Parallel integer relation
detection: techniques and applications,''
\textit{Math.~Comp.}\ \textbf{70}, 1719--1736 (2001).

\bibitem{apery1979}
R.~Ap\'ery, ``Irrationalit\'e de $\zeta(2)$ et $\zeta(3)$,''
\textit{Ast\'erisque}\ \textbf{61}, 11--13 (1979).

\end{thebibliography}

\end{document}
